\documentclass[11pt]{article}

\usepackage[margin=1in]{geometry}
\usepackage{amsmath, amssymb}
\usepackage{natbib}
\usepackage{hyperref}

\title{Clinical background}
\author{Tijn Jacobs}
\date{\today}

\begin{document}
\maketitle

\subsection{Clinical background and data sources}

Our analysis combines two data sources on the treatment of HIV-1 infection.
The randomized component is the AIDS Clinical Trials Group Study~175
(ACTG~175), a double-blind trial that enrolled $2{,}139$ HIV-1--infected
adults with CD4 cell counts between $200$ and $500$ cells/mm$^3$ and
randomized them to one of four nucleoside reverse-transcriptase inhibitor
regimens: zidovudine (ZDV) monotherapy, didanosine (ddI) monotherapy, ZDV
combined with ddI, or ZDV combined with zalcitabine (ddC)
\citep{hammer1996}. The observational component is drawn from the
Multicenter AIDS Cohort Study (MACS), a long-running prospective cohort of
men who have sex with men that has documented the natural and treated
history of HIV-1 infection since 1984 \citep{kaslow1987}. We restrict both
sources to a comparable pre-HAART calendar window and to participants
meeting the trial's eligibility criteria, and we define a binary treatment
contrast: ZDV monotherapy versus the pooled alternative of ddI monotherapy,
ZDV+ddI, or ZDV+ddC. The outcome is event-free survival, defined as the
time to the first of a $\geq 50\%$ decline in CD4 count, an AIDS-defining
event, or death.

The clinical evidence on this treatment contrast is well established.
ACTG~175 found that ZDV monotherapy was inferior to each of the other three
regimens, which were in turn comparable to one another: ddI monotherapy,
ZDV+ddI, and ZDV+ddC each delayed disease progression and prolonged
survival relative to ZDV alone \citep{hammer1996}. The Delta trial reached
the same conclusion, showing that combinations of ZDV with ddI or ddC were
superior to ZDV monotherapy \citep{delta1996}. The pharmacological
explanation is the rapid emergence of resistance under ZDV monotherapy:
HIV-1 reliably accumulates thymidine-analogue mutations in reverse
transcriptase within months of treatment, eroding any initial benefit
\citep{larder1989}. While zidovudine had earlier been shown to confer a
short-lived advantage over deferred treatment in symptom-free individuals
\citep{concorde1994}, this transient benefit reflects a comparison against
no treatment and does not extend to a comparison against ddI- or
combination-based regimens. Pooling the three superior arms into a single
comparator is therefore both clinically and statistically justified.

This contrast is well suited to studying the integration of randomized and
observational evidence precisely because the two sources are expected to
disagree. In observational practice, the choice of regimen was driven by
disease severity, immunologic trajectory, and prior treatment experience,
so that patients with poorer prognosis were preferentially escalated to
ddI-containing or combination therapy. This channeling induces confounding
by indication that operates opposite to the true treatment effect, and can
produce a spurious early survival advantage for ZDV monotherapy in the
observational cohort---together with an apparent attenuation of that
advantage over follow-up, as the mis-sorted high-risk patients experience
events and the susceptible pool is depleted \citep{hernan2010}. Neither
feature is present under randomization. The discrepancy between the
unconfounded trial estimate and the confounded observational estimate is
precisely the quantity that the confounding function in our framework is
designed to capture.

\begin{itemize}
  \item ACTG~175 was a double-blind randomized trial of $2{,}139$ HIV-1--infected
        adults with CD4 counts of $200$--$500$ cells/mm$^3$, randomized to ZDV
        monotherapy, ddI monotherapy, ZDV+ddI, or ZDV+ddC; the primary endpoint was
        event-free survival, i.e.\ time to a $\geq 50\%$ CD4 decline, an
        AIDS-defining event, or death \citep{hammer1996}.

  \item The trial's principal finding was that ZDV monotherapy was inferior to each
        of the other three regimens, which were statistically comparable to one
        another --- providing the clinical justification for collapsing them into a
        single comparator arm \citep{hammer1996}.

  \item The Delta trial independently corroborated this result, showing ZDV+ddI and
        ZDV+ddC superior to ZDV monotherapy \citep{delta1996}.

  \item The inferiority of ZDV monotherapy is mechanistically explained by the rapid
        emergence of thymidine-analogue resistance mutations in HIV-1 reverse
        transcriptase, which erode the initial treatment response within months of
        starting therapy \citep{larder1989}.

  \item The early, short-lived benefit of ZDV documented in Concorde was a comparison
        against \emph{deferred or no} treatment in symptom-free individuals; it does
        not imply any advantage of ZDV monotherapy over ddI- or combination-based
        regimens, and so does not apply to the present treatment contrast
        \citep{concorde1994}.

  \item MACS is a prospective observational cohort of men who have sex with men that
        has followed the natural and treated history of HIV-1 infection since 1984;
        regimen choice was made in routine clinical care rather than randomized
        \citep{kaslow1987}.

  \item In observational practice, regimen choice tracked disease severity, CD4
        trajectory, and prior treatment experience, so that poorer-prognosis patients
        were preferentially escalated to ddI-containing or combination therapy
        (confounding by indication).

  \item This channeling sorts high-risk patients into the comparator arm and can make
        ZDV monotherapy appear spuriously favorable in the observational data ---
        opposite in direction to the randomized result --- with an apparent
        attenuation of that advantage over follow-up as susceptible patients are
        depleted \citep{hernan2010}.

  \item The introduction of protease inhibitors and HAART in 1996 redefined standard
        of care; both data sources must be restricted to a comparable pre-HAART
        calendar window for the treatment definitions to remain meaningful and
        exchangeable.

  \item ACTG~175 was stratified by prior antiretroviral use; ZDV-experienced patients
        are more likely to harbor pre-existing resistance, making prior ARV exposure
        a clinically motivated candidate effect modifier for the treatment effect.

  \item In the observational cohort, regimens were switched as disease progressed; a
        baseline-fixed treatment definition therefore leaves later follow-up subject
        to switching-induced selection.

  \item Regimen toxicity profiles differ (ZDV: bone-marrow suppression; ddI:
        pancreatitis and peripheral neuropathy; ddC: peripheral neuropathy), which
        can contribute to early, regimen-specific attrition.
\end{itemize}



\subsection{The error distribution}

Two features of the application dictate the form of the error distribution: it
must be nonparametric, and it must be specific to the data source.

\paragraph{Why nonparametric.} The event-free survival endpoint is a composite,
defined as the first occurrence of a $\geq 50\%$ decline in CD4 count, an
AIDS-defining event, or death \citep{hammer1996}. The time to the first of
several biologically distinct events does not follow any standard parametric
family, and the log-normal, Weibull, and log-logistic errors implied by the
usual parametric accelerated failure time models cannot be expected to capture
its shape. Crucially, in this setting error misspecification is not innocuous. A
parametric accelerated failure time model fully specifies the likelihood, so an
incorrect error shape is accommodated by distorting the location parameters---the
baseline function $m_0$, the treatment effect $\tau$, and the confounding
function $c$. Because these location parameters are precisely the causal
estimands of interest, a nonparametric error is required to protect them from
shape misspecification, rather than merely to improve the fit of the survival
curve.

\paragraph{Why source-specific.} A single nonparametric error, shared between the
trial and the cohort, still imposes the assumption that the two samples have an
identical residual distribution. This homogeneity assumption is difficult to
defend here. First, the decomposition is a working model marginal over the
unmeasured confounders $U$; because treatment is selected on $U$ in the
observational study but randomized in the trial, the residual distribution of
the treated observational arm is reshaped relative to the trial, and this
reshaping affects higher moments that the confounding function $c(X)$, capturing
only the location, does not absorb. Second, the endpoint is ascertained
differently---on a protocol schedule with adjudication in the trial, and by
reconstruction from approximately six-monthly cohort visits in the observational
study, which coarsens event times. Third, the censoring and follow-up regimes
differ, the cohort being additionally subject to loss to follow-up and to
truncation at the onset of the HAART era. Fourth, the observational population
is broader and less selected for adherence than the trial population. These
mechanisms predict differences in the \emph{shape} of the residual
distribution---its skewness, tails, and modality---and not merely its scale, so
that a shared error with a source-specific variance would also be inadequate. At
the same time, the two samples concern the same disease, endpoint, and time
scale, so their error distributions are related rather than arbitrary. The
hierarchical Dirichlet process is the appropriate model for this situation: it
equips each source with its own error distribution while sharing a common pool
of mixture components, and it reduces to a single shared distribution when the
two samples are in fact homogeneous. The hierarchical specification therefore
relaxes the homogeneity assumption without requiring that the sources be shown
to differ in advance.


\textbf{The error distribution: the argument.}

\medskip
\noindent\emph{Part A --- the error must be nonparametric.}
\begin{itemize}
  \item \textbf{R1.} The endpoint is a composite---the first of a $\geq 50\%$
        CD4 decline, an AIDS-defining event, or death---and the time to the
        first of such heterogeneous events has no standard parametric form.
  \item \textbf{R2.} In an AFT model the likelihood is fully specified, so a
        misspecified error is absorbed by the location parameters
        $m_0,\tau,c$. Since these are the causal estimands, parametric
        misspecification biases the targets of inference.
\end{itemize}

\noindent\emph{Part B --- the error must be source-specific (HDP, not a single DP).}
A single DP assumes the trial and the cohort share one error distribution; this
homogeneity assumption fails for four reasons.
\begin{itemize}
  \item \textbf{R3 (structural).} The model is marginal over the unmeasured $U$;
        treatment is selected on $U$ in the cohort but randomized in the trial,
        reshaping the cohort's residual distribution. $c(X)$ corrects the
        location of this distortion, not its shape.
  \item \textbf{R4 (measurement).} The endpoint is adjudicated on a protocol
        schedule in the trial but reconstructed from six-monthly visits in the
        cohort, coarsening event times.
  \item \textbf{R5 (design).} Censoring differs: administrative in the trial;
        loss to follow-up and HAART-era truncation in the cohort.
  \item \textbf{R6 (population).} The cohort is broader and less
        adherence-selected than the trial.
\end{itemize}
Two further steps complete the argument.
\begin{itemize}
  \item \textbf{R7.} R3--R6 alter the \emph{shape} of the error---skew, tails,
        modality---not merely its scale; a shared DP with a source-specific
        variance is therefore also insufficient.
  \item \textbf{R8.} The two sources share disease, endpoint, and time scale, so
        independent DPs would discard common structure.
\end{itemize}

\noindent\textbf{Conclusion.} R1--R2 require a nonparametric error; R3--R8
require it to be source-specific with partial pooling. The hierarchical
Dirichlet process is the implied model---source-specific error distributions
over a shared atom pool---and it reduces to a single DP when homogeneity holds.




\vspace{1cm}

\noindent\textbf{Important questions.}
\begin{itemize}
  \item From when do we start counting time? For the RCT since
    randomisation; how about the RWD?
  \item Go into the selection of the RCT---restrict the RWD to a
    similar cohort?
  \item How can we use the WIHS data? Merge it with the MACS, treat it
    as a separate source, or use it as a validation set?
  \item Interval censored: currently the left interval point is
    registered as the event time.
\end{itemize}

\newpage


\bibliographystyle{apalike}
\bibliography{../references}

\end{document}
